
\section{The interaction surface}\label{sec:surface}

Houston's formulation~\cite{HOU16B} identifies full magic with three
same-difference square progressions whose centres themselves form an
arithmetic progression. We first retain the fixed difference $840$ of the
preceding construction, then allow the difference to vary and remove the
common scale of the roots. These are different arithmetic models. The
general setting of linear coordinate relations in $E^3$ is studied by Caro
and Garc\'ia-Fritz~\cite{CGF25}; their finiteness theorem is specialized in
Section~\ref{sec:finiteness}.

\begin{theorem}[The interaction surface]\label{main:surface}
\emph{(i)} After independent common subtractions $\tau_0,\tau_1,\tau_2$ on
the three difference-$840$ transversals of the fixed carrier, all rows and
columns remain equal and the two diagonals agree if and only if
$\tau_1+\tau_2=2\tau_0$. Requiring square entries gives
$x_i\in x\bigl(2\Esn(\Q)\bigr)$ and the interaction $x_1+x_2=2x_0$.
\emph{(ii)} With nonzero common difference $d$ allowed to vary, the
signed-root interaction model, modulo common scaling of the roots, has
the irreducible surface $Y^\circ$ defined below as its nonzero,
distinct-entry locus. Its map to the rational cone $2A^2=B^2+C^2$ has
generic degree five in squared coordinates and degree forty on the full
signed cover; the quintic reconstruction over the rational function field
of the cone has Galois group $S_5$.
For a prescribed difference $d_0$, recovering the roots additionally
requires a nonzero scale $L$ satisfying
\[
L^2=\frac{2d_0W}{RS}.
\]
The irreducibility assertion concerns $Y^\circ$, and the Galois-group
assertion concerns its quintic reconstruction, without an assertion about
this fixed-difference scale locus or the group of the full signed cover.
\end{theorem}

\subsection{The transversal equation}

For arbitrary $x_0,x_1,x_2$ consider
\[
\begin{pmatrix}
x_1 & x_2+840 & x_0-840\\
x_2-840 & x_0 & x_1+840\\
x_0+840 & x_1-840 & x_2
\end{pmatrix}.
\]
Every row, every column and the principal diagonal sums to $x_0+x_1+x_2$,
while the antidiagonal sums to $3x_0$; all eight sums agree if and only if
\[
x_1+x_2=2x_0 .
\]
This proves Theorem~\ref{main:surface}(i): imposing that each $x_i$ be the
$x$-coordinate of a point of $2\Esn(\Q)$ --- the arithmetic condition of
Lemma~\ref{lem:doubling} --- makes the fully magic completion of the
transversal frame exactly one linear equation on $\Esn^{3}$, hence a
surface. The diagonal $(P,P,P)$ is not isolated, and the membership
condition $P_i\in2\Esn(\Q)$ is arithmetic rather than a further geometric
equation over an algebraic closure.

For later use we record the Mordell--Weil data. The minimal model of
$\Esn$ is recorded as LMFDB \texttt{705600.vn3}~\cite{LMFDB},
$Y^2=X^3-44100X$, with
rank $2$, full rational $2$-torsion and free generators $(-84,1764)$,
$(294,3528)$; transporting by $(x,y)=(4X,8Y)$ and adjusting by the
two-torsion point $T=(-840,0)$ gives the free basis
\[
Q=(1960,78400),\qquad R=(1176,28224)
\]
used throughout, with
\[
2Q=(841,-1189),\qquad 2R=(1369,-39997).
\]
Since doubling kills the rational two-torsion, $2Q$ and $2R$ generate the
free part of $2\Esn(\Q)$.

\subsection{Normalizing the model with variable difference}

Let $(r_i^2,s_i^2,t_i^2)$, $i=0,1,2$, be three progressions of rational
squares with a common difference $d\ne0$, and retain a choice of their
root signs. Put
\[
u_i=t_i-r_i,\qquad v_i=t_i+r_i.
\]
The progression equations and the full-magic interaction become
\[
u_iv_i=2d,\qquad u_i^2+v_i^2=4s_i^2,\qquad
s_1^2+s_2^2=2s_0^2.
\]
Conversely these equations recover the roots
$r_i=(v_i-u_i)/2$, $t_i=(v_i+u_i)/2$. Since $d\ne0$, every $u_i,v_i$
is nonzero. Define the scalar coordinates
\[
\begin{gathered}
L=u_0,\qquad R=\frac{u_1}{u_0},\qquad
S=\frac{u_2}{u_0},\qquad W=\frac{u_1u_2}{2d},\\
(A,B,C)=\frac{u_1u_2}{d\,u_0}(s_0,s_1,s_2).
\end{gathered}
\]
Here the scalar $R$ is unrelated to the elliptic point denoted by $R$
above. Direct substitution gives the affine variety $Y$:
\[
\begin{aligned}
F&=2R^2S^2-(1+W^2)(R^2+S^2)+2W^2=0,\\
A^2&=W^2+R^2S^2, & B^2&=W^2R^2+S^2, & C^2&=W^2S^2+R^2.
\end{aligned}
\]
The inverse map, valid for $RSWL\ne0$, is
\[
\begin{aligned}
d&=\frac{RSL^2}{2W},&
(u_0,u_1,u_2)&=(L,RL,SL),\\
(v_0,v_1,v_2)&=\frac L W(RS,S,R),&
(s_0,s_1,s_2)&=\frac L{2W}(A,B,C).
\end{aligned}
\]
In particular, the three triples of roots are
\begin{equation}\label{eq:surface-root-map}
\begin{aligned}
(r_0,s_0,t_0)&=\frac L{2W}(RS-W,A,RS+W),\\
(r_1,s_1,t_1)&=\frac L{2W}(S-RW,B,S+RW),\\
(r_2,s_2,t_2)&=\frac L{2W}(R-SW,C,R+SW).
\end{aligned}
\end{equation}
The radicand equations give the two differences in each progression as
$d=RSL^2/(2W)$, and $F=0$ gives $2A^2=B^2+C^2$, hence the interaction.
Substituting in either direction recovers every coordinate. Thus the
nonzero signed-root model with variable $d$ is isomorphic to the
corresponding locus of $Y\times\mathbb G_m$, with $L$ the
$\mathbb G_m$ coordinate. Multiplying all roots by $q\ne0$ multiplies
$L$ by $q$ and $d$ by $q^2$, leaving $R,S,W,A,B,C$ unchanged.

For a fixed difference $d_0$, this inverse map instead requires
\begin{equation}\label{eq:surface-fixed-scale}
L^2=\frac{2d_0W}{RS},\qquad
(x_0,x_1,x_2)=\frac{d_0}{2WRS}(A^2,B^2,C^2).
\end{equation}
When $L$ is rational, \eqref{eq:surface-root-map} supplies rational roots
of $x_i-d_0,x_i,x_i+d_0$, the square-progression data associated with
$x\bigl(2E_{d_0}(\Q)\bigr)$. Conversely, these data give the coordinates
and scale above. For $d_0=840$ the additional equation is
$L^2=1680W/(RS)$. In particular, an arbitrary rational point of $Y$
does not supply rational roots at difference $840$: the quantity on the
right must be a rational square. The normalization retains the
squareclass $[d]=[RS/(2W)]$ and does not identify all differences with one
fixed value.

\subsection{The nonzero distinct-entry locus and saturation}

The excluded divisors can be read directly from the progression data.
Set
\[
(\eta_0,\eta_1,\eta_2)=
\left(\frac{2W}{RS},\frac{2WR}{S},\frac{2WS}{R}\right),
\qquad f(\eta)=\frac{\eta^2+4}{4\eta}.
\]
These satisfy $\eta_i=u_i^2/d$ and $s_i^2/d=f(\eta_i)$, so the nine
entries divided by $d$ are $f(\eta_i)+\varepsilon$ for
$i=0,1,2$ and $\varepsilon\in\{-1,0,1\}$. Define $Y^\circ$ by removing
$RSWABC=0$, the numerators of the nine equations
$f(\eta_i)+\varepsilon=0$, and those of the $27$ equations
\[
f(\eta_i)+\varepsilon=f(\eta_j)+\varepsilon',
\qquad i<j,\quad \varepsilon,\varepsilon'\in\{-1,0,1\}.
\]
Within a progression the entries are already distinct because $d\ne0$.
The $27$ collision numerators have $30$ distinct normalized factors.
The coordinate maps preserve these exclusions.
Positivity remains a real semialgebraic condition: for positive increasing
roots one takes $L,R,S,W,A,B,C>0$ and
$RS>W$, $S>RW$, $R>SW$. It is not an additional polynomial saturation.

\begin{proof}[Proof of Theorem~\ref{main:surface}(ii), saturation part]
As a polynomial in $S$,
\[
F=(2R^2-W^2-1)S^2+2W^2-(1+W^2)R^2,
\]
whose leading coefficient and constant term are coprime; the coefficient
$2R^2-W^2-1$ is squarefree and occurs to odd order in the ratio defining
$S^2$, so that ratio is not a square in $\Q(R,W)$ and $F$ is irreducible.
Over the function field of $F=0$ put $p_A=W^2+R^2S^2$, $p_B=W^2R^2+S^2$,
$p_C=W^2S^2+R^2$. In squared coordinates $r=R^2$, $s=S^2$, $w=W^2$, setting
$p_A=0$ gives $F(r,s,-rs)=(r+s)(rs-1)$, and on the reduced component
$RS=1$, $W^2=-1$ the function $p_A$ has generic valuation one while $p_B$
and $p_C$ are generically nonzero; the case $p_B=0$ reduces to
$-2r^2w+rw^2-r+2w=0$, linear in $p_B$ at the generic point of each reduced
component with $p_A$ and $p_C=2p_A$ generically nonzero; the $p_C$ case is
symmetric. Hence every nonempty product of the three radicands has odd
valuation somewhere and is not a square, so adjoining $A,B,C$ gives a
degree-eight extension with integral coordinate ring: the four equations
define a single prime component of dimension two. None of the declared
zero/collision polynomials vanishes on it, so saturating the prime ideal by
their product changes nothing, and $Y^{\circ}$ is one irreducible surface.
\end{proof}

\subsection{The square-progression cone and quintic reconstruction}

The four equations imply the exact identity
\[
2A^2=B^2+C^2 :
\]
the three auxiliary roots themselves form a square progression, and the
normalized interaction maps $Y$ dominantly to this classical cone. In squared
coordinates $a=rs+w$, $b=rw+s$, $c=sw+r$ (so $2a=b+c$), reconstruction away
from $r^2=1$ is
\[
s=\frac{b-ar}{1-r^2},\qquad w=\frac{a-br}{1-r^2},
\]
where $r$ satisfies the quintic
\[
Q(r)=c(1-r^2)^2-(b-ar)(a-br)-r(1-r^2)^2=0 .
\]

\begin{proof}[Proof of Theorem~\ref{main:surface}(ii), cone part]
The quintic has leading coefficient $-1$ and the displayed inverse
reconstructs $s,w$ from each generic root, so the squared-coordinate map to
the cone is generically finite of degree five; each squared triple has
eight sign choices of $(R,S,W)$, giving degree $5\cdot8=40$ on the full
signed cover. For the Galois group over the rational function field of the
cone, specialize at the rational cone point
$(A,B,C)=(5,1,7)$: the reconstruction polynomial
\[
-r^5+49r^4+2r^3-123r^2+625r+24
\]
is irreducible modulo $7$, factors modulo $29$ with degree pattern
$[2,1,1,1]$, and has nonzero discriminant $-24364389070416096000$, so the
specialization is separable. Irreducibility supplies a transitive
degree-five group containing a five-cycle; the mod-$29$ pattern supplies a
transposition; among transitive subgroups of $S_5$ this forces the full
$S_5$, and the generic Galois group contains the specialized one. Since
the generic group is itself a subgroup of $S_5$, it is exactly $S_5$.
\end{proof}

\begin{remark}[Scope of the quintic obstruction]
The generic $S_5$ Galois group proves that the cone map has no generic rational
section and that its quintic cannot be resolved by a tower of quadratic
radicals. It does \emph{not} prove that $Y$ admits no other useful rational
map, no height-decreasing correspondence, and no further structure; those
remain open research directions, not consequences.
These conclusions concern the normalized model with variable difference.
No irreducibility or Galois-group conclusion for the fixed-$d_0$ locus
defined by \eqref{eq:surface-fixed-scale} is inferred here.
\end{remark}
